\typeout{=== RC_Pass START ===}
\subsection{RC Pass Filter}

\subsubsection{Overview}
\paragraph
Constructed from a resistor and capacitor in series. It
is a single pole filter. An RC Pass Filter can be broken
down into two types: Low Pass and High Pass.
\begin {enumerate}
  \item \textbf{Low Pass Filter:} Allows low frequency
    signals to pass while attenuating high frequency
    signals.
  \item \textbf{High Pass Filter:} Allows high frequency
    signals to pass while attenuating low frequency
    signals.
\end{enumerate}
\par
\subsubsection{Schematic}
\begin{circuitfig}
  % Paths, nodes and wires:
  \draw (4.25, 6.5) -- (5.25, 6.5);
  \node[shape=rectangle, minimum width=0.715cm, minimum height=0.465cm](N1) at (5.625, 6.5){} node[anchor=center] at (N1.text){$V_{out}$};
  \draw (4.25, 6.5) to[capacitor, l={$C_1$}] (4.25, 4.75);
  \draw (4.25, 6.5) to[european resistor, l_={$G_1$}, label distance=0.03cm] (4.25, 8.25);
  \draw (2.75, 8.25) to[sinusoidal voltage source, l={$V_1$}] (2.75, 6.75);
  \draw (2.75, 4.75) to[sinusoidal voltage source, l_={$V_2$}] (2.75, 6.25);
  \node[ground, rotate=-90] at (2.75, 6.75){};
  \node[ground, rotate=-90] at (2.75, 6.25){};
  \draw (2.75, 4.75) -- (4.25, 4.75);
  \draw (2.75, 8.25) -- (4.25, 8.25);
  \node[shape=rectangle, minimum width=2.965cm, minimum height=0.465cm] at (3.5, 4.25){} node[anchor=north, align=center, text width=2.577cm, inner sep=6pt] at (3.5, 4.5){Full Schematic};
  \draw (8.625, 6.5) -- (9.625, 6.5);
  \node[shape=rectangle, minimum width=0.715cm, minimum height=0.465cm](N2) at (10, 6.5){} node[anchor=center] at (N2.text){$V_{out}$};
  \draw (8.625, 6.5) to[capacitor, l={$C_1$}] (8.625, 4.75);
  \draw (8.625, 6.5) to[european resistor, l_={$G_1$}, label distance=0.03cm] (8.625, 8.25);
  \draw (7.125, 8.25) to[sinusoidal voltage source, l={$V_1$}] (7.125, 6.75);
  \node[ground, rotate=-90] at (7.125, 6.75){};
  \node[ground, rotate=-90] at (7.125, 4.75){};
  \draw (7.125, 4.75) -- (8.625, 4.75);
  \draw (7.125, 8.25) -- (8.625, 8.25);
  \node[shape=rectangle, minimum width=2.965cm, minimum height=0.465cm] at (7.875, 4.25){} node[anchor=north, align=center, text width=2.577cm, inner sep=6pt] at (7.875, 4.5){Low Pass};
  \draw (13, 6.5) -- (14, 6.5);
  \node[shape=rectangle, minimum width=0.715cm, minimum height=0.465cm](N3) at (14.375, 6.5){} node[anchor=center] at (N3.text){$V_{out}$};
  \draw (13, 6.5) to[capacitor, l={$C_1$}] (13, 4.75);
  \draw (13, 6.5) to[european resistor, l_={$G_1$}, label distance=0.03cm] (13, 8.25);
  \draw (11.5, 4.75) to[sinusoidal voltage source, l_={$V_2$}] (11.5, 6.25);
  \node[ground, rotate=-90] at (11.5, 8.25){};
  \node[ground, rotate=-90] at (11.5, 6.25){};
  \draw (11.5, 4.75) -- (13, 4.75);
  \draw (11.5, 8.25) -- (13, 8.25);
  \node[shape=rectangle, minimum width=2.965cm, minimum height=0.465cm] at (12.25, 4.25){} node[anchor=north, align=center, text width=2.577cm, inner sep=6pt] at (12.25, 4.5){High Pass};
\end{circuitfig}

\subsubsection{Transfer Function}
Applying KCL at $V_{out}$ on Full Schematic, we get:
\[
  (V_{out} - V_{1}){G_1} + (V_{out} - V_{2}){sC_1} = 0
\]
\paragraph{Superposition:}
\begin{enumerate}
  \item Due to \(V_{1}\) alone (\(V_{2} = 0V\)):
    This is the low-pass characteristic of the RC Pass Filter.
    \begin{align*}
      (V_{out} - V_{1}){G_1} + (V_{out} - 0){sC_1} &= 0 \\
      V_{out}G_1 - V_{1}G_1 + V_{out}sC_1 &= 0 \\
      V_{out}(G_1 + sC_1) - V_{1}G_1 &= 0 \\
      V_{out}(G_1 + sC_1) &= V_{1}G_1 \\
      H_1(s) = \frac{N_1(s)}{D_1(s)} = \frac{V_{out_1}}{V_{1}} &= \frac{G_1}{G_1 + sC_1}
    \end{align*}
  \item Due to \(V_{2}\) alone (\(V_{1} = 0V\)):
    This is the high-pass characteristic of the RC Pass Filter.
    \begin{align*}
      (V_{out} - 0){G_1} + (V_{out} - V_{2}){sC_1} &= 0 \\
      V_{out}G_1 + V_{out}sC_1 - V_{2}sC_1 &= 0 \\
      V_{out}(G_1 + sC_1) - V_{2}sC_1 &= 0 \\
      V_{out}(G_1 + sC_1) &= V_{2}sC_1 \\
      H_2(s) = \frac{N_2(s)}{D_2(s)} = \frac{V_{out_2}}{V_{2}} &= \frac{sC_1}{G_1 + sC_1}
    \end{align*}
\end{enumerate}

\paragraph{Total Output:}
The full output of the RC Pass Filter is the sum of the
low-pass and high-pass characteristics.
\begin{align*}
  V_{out} &= V_{out_1} + V_{out_2} \\
  V_{out} &= \underbrace{\frac{V_{1}G_1}{G_1 + sC_1}}_{\text{Low Pass Character}} + \underbrace{\frac{V_{2}sC_1}{G_1 + sC_1}}_{\text{High Pass Character}}
\end{align*}

\subsubsection{Analysis}
\begin{enumerate}
  \item \textbf{Poles (\(D(s) = 0\)):} \\
    For both low pass \(H_1(s)\) and high pass \(H_2(s)\):
      \begin{align*}
        D(s) \quad\rightarrow\quad G_1 + sC_1 = 0 \\
        s_p &= -\frac{G_1}{C_1} \\
      \end{align*}
    Negative real pole indicates a stable first-order system.
    Frequency corresponding to this pole:
    \begin{align*}
      s &= -2 \pi f_p \\
      -2 \pi f_p &= -\frac{G_1}{C_1} \\
      f_p = \frac{G_1}{2 \pi C_1} = \frac{1}{2 \pi R_1 C_1}
    \end{align*}
  \item \textbf{Zeroes (\(N(s) = 0\)):} \\
    For low pass $H_1(s)$:
      \begin{align*}
        N_1(s) &= 0 \\
        G_1 &\neq 0
      \end{align*}
    No zeroes for \(H_1(s)\). \\
    For high pass $H_2(s)$:
      \begin{align*}
        N_2(s) &= 0 \\
        sC_1 &= 0 \\
        s_z &= 0 \quad \text{[let \(s = s_z\)]}\\
        f_z &= \frac{|s_z|}{2\pi} = 0 Hz
      \end{align*}
    Zero at origin for \(H_2(s)\).
  \item \textbf{DC Gain (\(H(0)\)):} \\
    For low pass \(H_1(s)\):
      \begin{align*}
        H_1(0) &= \frac{G_1}{G_1 + 0 \cdot C_1} = 1 \\
        A_{DC_1} = 1 = 0 dB
      \end{align*}
    Unity DC gain for \(H_1(s)\). \\
    For high pass \(H_2(s)\):
      \begin{align*}
        H_2(0) &= \frac{0 \cdot C_1}{G_1 + 0 \cdot C_1} = 0 \\
        A_{DC_2} = 0 = -\infty dB
      \end{align*}
    Zero DC gain for \(H_2(s)\).
  \item \textbf{Gain-Bandwidth Product (\(GBW\)):}
    For low pass \(H_1(s)\):
    \begin{align*}
      \text{GBW} &= |A_{DC_1} \times f_p| \\
                 &= |1 \times \frac{1}{2\pi R_1C_1}| \\
                 &= \frac{1}{2\pi R_1C_1}
    \end{align*}
    For high pass \(H_2(s)\):
    \begin{align*}
      \text{GBW} &= |A_{DC_2} \times f_p| \\
                 &= |0 \times \frac{1}{2\pi R_1C_1}| \\
                 &= 0
    \end{align*}
    The high pass filter has no gain at DC, so the GBW which
    is a quantity measured from DC to the cutoff frequency is
    zero.
  \item \textbf{Impedance:}
    Both low pass and high pass filters have the same input
    and output impedance.
    Substituting \(G_1 = \frac{1}{R_1}\) and \(s = j \omega\):
    \paragraph{At Input:}
      \begin{align*}
        Z_{in} = R_1 + \frac{1}{sC_1} \\
        \text{Substituting } s = j \omega &: \\
        Z_{in} = R_1 - j \frac{1}{\omega C_1}
      \end{align*}
    \paragraph{At Output:}
      \begin{align*}
        Z_{out} &= R1 \parallel \frac{1}{sC_1} \\
        Z_{out} &= \frac{R1 \cdot \frac{1}{sC_1}}{R1 + \frac{1}{sC_1}} \\
        Z_{out} &= \frac{\frac{R_1}{sC_1}}{\frac{sC_1 R_1 + 1}{sC_1}} \\
        Z_{out} &= \frac{R_1}{sC_1 R_1 + 1} \\
        \text{Substituting } s = j \omega &: \\
        Z_{out} &= \frac{R_1}{1 + j \omega R_1 C_1} \cdot \frac{1 - j \omega R_1 C_1}{1 - j \omega R_1 C_1} \\
        Z_{out} &= \frac{R_1 - j \omega R_1^2 C_1}{1 + (\omega R_1 C_1)^2} \\
        Z_{out} = \frac{R_1}{1 + (\omega R_1 C_1)^2} - j \frac{\omega R_1^2 C_1}{1 + (\omega R_1 C_1)^2}
      \end{align*}
\end{enumerate}

\subsubsection{Question}
Design a RC Low Pass Filter with a cutoff frequency of
\(f_p = 45kHz\) using a capacitor value of \(C_1 = 1\mu F\).
Also compute the input and output impedance at
\(f = 10kHz\).
\paragraph{Solution:}
\begin{enumerate}
  \item Calculating \(R_1\):
    \begin{align*}
      f_p &= \frac{1}{2 \pi R_1 C_1} \\
      R_1 &= \frac{1}{2 \pi f_p C_1} \\
      &= \frac{1}{2 \pi (45kHz)(1\mu F)} \\
      R_1 &\approx 3.54 \Omega
    \end{align*}
  \item Calculating \(Z_{in}\) and \(Z_{out}\) at \(f = 10kHz\):
    \begin{align*}
      \omega &= 2 \pi f = 2 \pi (10kHz) = 62,832 \text{ rad/s} \\
      Z_{in} &= R_1 - j \frac{1}{\omega C_1} \\
             &= 3.54 \Omega - j \frac{1}{(62,832)(1\mu F)} \\
             &\approx 3.54 \Omega - j 15.9 \Omega \\
    \end{align*}
    \begin{align*}
            Z_{out} &= \frac{R_1}{1 + (\omega R_1 C_1)^2} - j \frac{\omega R_1^2 C_1}{1 + (\omega R_1 C_1)^2} \\
              &= \frac{3.54 \Omega}{1 + (62,832 \cdot 3.54 \Omega \cdot 1\mu F)^2} - j \frac{(62,832)(3.54 \Omega)^2 (1\mu F)}{1 + (62,832 \cdot 3.54 \Omega \cdot 1\mu F)^2} \\
              &= \frac{3.54 \Omega}{1 + (0.222)^2} - j \frac{(62,832)(12.52 \Omega^2)(1\mu F)}{1 + (0.222)^2} \\
              &\approx \frac{3.54 \Omega}{1.049} - j \frac{0.786 \Omega}{1.049} \\
              &\approx 3.37 \Omega - j 0.75 \Omega
    \end{align*}
\end{enumerate}
Design a high pass filter with a cutoff frequency of
\(f_p = 1kHz\) using a capacitor value of \(C_1 = 1\mu F\).
Also compute the input and output impedance at
\(f = 75Hz\).
\paragraph{Solution:}
\begin{enumerate}
  \item Calculating \(R_1\):
    \begin{align*}
      f_p &= \frac{1}{2 \pi R_1 C_1} \\
      R_1 &= \frac{1}{2 \pi f_p C_1} \\
          &= \frac{1}{2 \pi (1kHz)(1\mu F)} \\
          &= 159.15 \Omega \\
      R_1 &\approx 159 \Omega
    \end{align*}
  \item Calculating \(Z_{in}\) and \(Z_{out}\) at \(f = 75Hz\):
    \begin{align*}
      \omega &= 2 \pi f = 2 \pi (75Hz) = 471.24 \text{ rad/s} \\
      Z_{in} &= R_1 - j \frac{1}{\omega C_1} \\
             &= 159 \Omega - j \frac{1}{(471.24)(1\mu F)} \\
             &\approx 159 \Omega - j 2123.72 \Omega \\
    \end{align*}
    \begin{align*}
          Z_{out} &= \frac{R_1}{1 + (\omega R_1 C_1)^2} - j \frac{\omega R_1^2 C_1}{1 + (\omega R_1 C_1)^2} \\
              &= \frac{159 \Omega}{1 + (471.24 \cdot 159 \Omega \cdot 1\mu F)^2} - j \frac{(471.24)(159 \Omega)^2 (1\mu F)}{1 + (471.24 \cdot 159 \Omega \cdot 1\mu F)^2} \\
              &\approx 153.96 \Omega - j 11.52 \Omega
    \end{align*}
\end{enumerate}
